\documentclass[tom-ari]{subfile}
\begin{document}
	
	\chapter{1 Major}
	
	\section{Intro}
	
	12 Jan 2020: Reverted to what we decided by the end of SF. This reflects playing \he1--\sp1 as natural rather than a relay.
	
	\section{Response summary} 

	Bids where they differ will be noted, but for the most part they retain the same meaning for both \he1 and \sp1.	
	
	\begin{bidtable}{\orauction{1M}}
		\sp1 & ``Standard'', 4+ cards forcing 1 round. \\
		1NT & Semi-Forcing NT \\
		\cl2 & 1 of 3 hand types: LR+, GF \clubsuit, GF balanced. Respond as per TaJ. \\ 
		\di2 & Natural GF. Allowed to bid either \cl2 or \di2 with 5332 no fit.\\
		\he2 & (After \sp1) Natural GF\\
		\sp2 & (After \he1) Natural, Inv \rem{Was mixed under the artificial \sp1 world, but changing to inv to match other jumps.} \\
		2NT & 5-5 minors, Mixed+.  Open to other meanings if useful. Strength open to discussion. \\
		3NT & Bal with M, choice of games.  No slam interest. \\
		R & Simple raise, less than inv. \\
		3x & Natural, Inv \\
		3M & Mixed \\
	\end{bidtable}

	\section{Semi-Forcing NT}

	2 minor bids are generally natural 4+ cards. Can be short if no other bid appeals and too strong to pass. (i.e., \di2 on \exactshape{5332}). Rebidding the other major is always natural; non-reverse can be any strength, reverse shows \shape{56xx} non-min, like a jump shift. 
	
	3 of a new suit is a 5--5 jump shift, non-forcing but maximum. 6--4 jump shifts are handled via 2NT. Over that, \cl3 is pass or correct, \di3 asks LMH (GF).
	
	Simple rebids and jump rebids in the major are both standard, 6+ NF.
		
	\section[2C TaJ]{\cl2 TaJ}
	
	This is the primary relay step.  It is inheriting the response structure from TaJ, but doesn't guarantee a fit the way TaJ does.  It is intended to be general purpose, which makes the bid mostly fall into 1 of 3 hand types:  a classic TaJ hand (LR+), any GF bal or other that feels like relaying, any GF in clubs.  The club hand can be shown by breaking relay and bidding \cl3 or \textbf{\textit{or higher. (i.e., \di3 would be primary clubs and secondary diamonds.)}} The only danger is that partner bids an \"Uber step 2NT or higher, but we are likely in a slam position so it may not matter much.
	
	\reversemarginpar
	\marginpar{\includegraphics[width=.8\marginparwidth]{new.png}\hfill} 
	
	The responses are as per TaJ, shifted down.  1st step is any min, 2nd step is any medium ``I would accept a LR'' type hand, steps above that are \"Uber good and break into TaJ steps.
	
	Over the more common +1/+2 responses, retreating to 2M is the weakest action, showing the LR type hand.  In the case of +2, this becomes the ``Are you sure?'' bid instead of raising to the 3 level.
	
	The cheapest bid that isn't 2M is the TaJ relay, as per 2NT in the old system.  Once we have relayed fully, we can advance into CIRKLE/SQUARE. 3NT is never an asking bid in relay auctions, it is always an attempt to play there. As in other auctions, when space is tight for CIRKLE then \cl4 becomes a puppet for any remaining CIRKLEs that can't be bid directly, \di4 is a puppet to sign off anywhere, and direct game bids above 3NT are NF but encouraging.

	A relay-then-3NT auction loses it's previous meaning of a Choice of Games, so the direct 1M--3NT responses have become that. The new meaning of relay-then-3NT is simply no slam interest opposite partner's strength, NF.
	
	\textbf{\textit{Added based on discussion with Jenni and Bryan:  2NT after the response (when not the relay) is a Woolsey style game try, asking for the lowest suit that Opener would accept a help suit GT.  The only time 2NT is a relay is \sp1--\cl2--\he2, so we should be in a more-or-less GF auction anyway, so a Woolsey GT isn't needed.}}
	
	\reversemarginpar
	\marginpar{\includegraphics[width=.8\marginparwidth]{new.png}\hfill} 
	
	\rem{Does it make sense to have a bail option to go straight to CIRKLE rather than needing to fully TaJ out? I can imagine, for example, hearing that partner is 54xx and not caring about the details, or even uncovering the 4 card suit and then bailing. Sample auction might be something like \sp1--\cl2--\di2--\he2--\sp2--2NT--\cl3--\he3.}
	
	\section[2D Nat GF]{\di2 Nat GF}
	
	GF with 5+ diamonds. You may choose to relay instead of showing diamonds, this is up to the responder. 2NT is the default response with nothing descriptive to say.
	
	\rem{Some people I play with like inverting 2M an 2NT here, making 2NT show 6+M and 2M be waiting. Thoughts?}
	
	\section[1S--2H]{\sp1--\he2}
	
	GF with 5+ hearts. You may choose to relay instead of showing diamonds, this is up to the responder.
	
	\rem{Removing old proposal about \sp1--\he2, instead going to have section about new methods.}
			
	\section{2NT}
	
	Both minors?  Really a place holder agreement, I'm open to anything. Mixed+? Other meanings might be better still, but I'm currently at a loss for finding a ``good'' use for this bid. Could be CIRKLE for M I suppose.
	
	\section{3NT}
	
	Choice of games, promises 3 card support for M.
	
	\section{Jump Shift}
	
	Natural, game invitational. 3M rebid by opener NF.
	
	\section{Double Jump Shift}
	
	Void Splinter. Singleton splinters start with TaJ. (\cl2)
	
	\chapter{1 Major - Proposal}

	\Huge{\color{red}ALPHA} \normalsize
	
	\section{Intro}
	
	This section is new as of 22 Aug 2020.  It may resemble what Ari and I tried to do in SF, but there are important differences.  Jenni has agreed to try it.
	
	This spouts from a few different things:
	
	\begin{itemize}
		\item Revisiting step based responses.
		\item Jenni/Bryan finding the 2C either/or TaJ unworkable
		\item Further examination of Moss/Grue methods
		\item Much soul searching
	\end{itemize}
		
	The idea is to have a step based responses as we did before, with keeping a separate TaJ-as-raise-only step.  The result is a mish-mosh of ideas we have explored or not.  I think this combined should work okay.
	
	\section{Response Summary}
	
	\begin{bidtable}{\orauction{1M}}
		+1 & Punt.  Like (Semi)Forcing NT.  Xfer responses to \sp1. \\
		+2 & GF, art. Typically denies 3+ card support\\
		+3 & TaJ, Limit+ \\
		+4 & Constructive to bad Invite in other major.  5+ cards, won't be 6+ in length and invitational.  (JS instead)  \\
		2M & Simple Raise \\
		JS & Nat Inv \\
		2NT & Mixed Raise - I talked with Jenni, she liked this to free up 1M--3M as pure weak. Also has symmetry with over opp's Dbl \\
		3M & Weak \\
		3NT & ??? \\
		DJS & Void Spl \\
	\end{bidtable}

	\section{+1: Punt}
	
	The +1 bid is a punt to keep the auction alive, much like a forcing or semi-forcing NT.  \sp1--1NT is non-forcing and frequently passed when balanced.  (We don't need to worry about the 3 card LR.)
	
	\subsection{\he1--\sp1}
	
	\sp1 by responder does not deny spades, although it is unlikely to hold 5+ \sss~and a decent hand.  (See +4)
	
	Over this punt, we play transfer style responses.  The 1NT rebid is either \shape{5332} or 4+ \ccc.
	
	\begin{bidtable}{\orauction{1h,1s}}
		1NT & 4+ \ccc~or balanced, NF \\
		\cl2 & 4+ \ddd \\
		\di2 & 6+ \hhh \\
		\he2 & Flannery; 4+ \sss \\
		\sp2 & 5-6 Nat 13+ \\
		2NT & ??? \\ 
		& \rem{With \di2 being (mostly) forcing, a good 6-4 hand can start there and then bid the 4 bagger.}\\
		\cl3/\ddd & 5-5 Nat 13+ \\
		\he3 & 6+ Max \\
	\end{bidtable}

	After Opener's rebid, no special methods.  Bids which sound like sign off are sign off.
	
	\subsection{\sp1--1NT}
	
	1NT is NF, essentially a classic Semi-Forcing NT.  Pass is encouraged.  Typically bidding over 1NT is 4+ natural, or a non-minimum.  \rem{2NT is a good 6--4?}
	
	\section{+2: GF}
	
	This is a \textit{non-relay} GF.  It can cover a wide range of GF hands, but it does not put us into a relay context automatically.  This allows for a (mostly) natural response structure by Opener, with a few tweaks (such as using 2NT to show extra length.)
	
	Since the GF denies a primary fit, we can use retreating to 2M when available as a relay.  Otherwise most follow up sequences by responder are natural.
	
	\subsection{\he1--1NT}
	
	Very natural here.  Plenty of room to show hand type and preserve immediate 5-5 jumps.  JS just show shape, not extras.  \he3 is normal 2/1 style jump, setting trump.  Instead of being limited to RKC or Qbids, we can immediately CIRKLE. (The last part is untested, but seems okay.) This makes +1 CIRKLE, +2 NS, +3/4 nat Serious Q.
	
	\begin{bidtable}{\orauction{1h,1n}}
		\cl2 & Exactly 4 \ccc, can have long \hhh \\
		\di2 & Exactly 4 \ddd, can have long \hhh \\
		\he2 & Exactly 4 \sss, can have long \hhh; pushed down for space reasons \\
		\sp2 & Any 5332 \\
		2NT & 6+ \hhh, non-solid \\
		3x & 5-5, any strength  (\sp3 is 5--6) \\
		\he3 & 0-1 loser suit, sets trumps. \sp3 over this is CIRKLE for \hhh, 3NT Non-S, 4m Serious \\
	\end{bidtable}

	\subsection{\sp1--\cl2}
	
	This is very similar to the old structure of \cl2, with some small modifications. \di2 is no longer a pure waiting bid but actively showing a 4 card minor. \sp2 shows the 5332, 2NT shows extra \sss.
	
	\begin{bidtable}{\orauction{1s,2c}}
		\di2 & Exactly 4 in either minor, can have long \sss \\
		\he2 & Exactly 4 \hhh, can have long \sss \\
		\sp2 & Any 5332, 2NT asks LMH. \\
		2NT & 6+ \sss, non-solid \\
		3x & 5-5, any strength \\
		\sp3 & 0-1 loser suit, sets trumps. 3NT over this is CIRKLE for \sss, \cl4 Non-S, 4Red Serious \\
	\end{bidtable}

	\subsection{Follow Ups}
	
	2M is used as a shape relay.  If the suit is known below 2M (\he1--1NT--2x or \sp1--\cl2--\he2) then you jump directly into shapes in a length count up:  5, 6, 7+ with zooming. Relay for shortness (or zoom) VLH(N/B). (None only for 5, both only for 7+.  6-4 must have shortness but cannot have both).
	
	If the suit is not known (\sp1--\cl2--\di2), step 1 is \ccc, then relay for 567.  Steps 2 and above are \ddd with the 567 count.
	
	Over \he1--1NT--\he2 (showing \sss), \sp2 becomes the shape relay for 567.  Note this leaves us 3NT as the 4522 step!
	
	Over \sp2 showing \shape{5332}, 2NT asks for the doubleton LMH.  
	
	Once shape is known, if Responder can bid \cl4 then we are in a CIRKLE situation.  3NT is always a sign off bid; 3 level bids are ABCD CIRKLE as available, with \cl4 puppets to \di4 when 2+ CIRKLE bids remain. (If only 1, \cl4 is that CIRKLE).  \di4 by Responder is a puppet to \he4 for sign off.  Note that all of this is fairly typical CIRKLE, just spelled out here for clarity.
	
	It is also quite likely that Responder will \textit{not} want to relay; if they choose any other actions than the ones listed above, they are natural in context.  2NT waiting, suits being natural, and so on.
	
	\section{+3: TaJ}
	
	3+ support, Limit Raise or better.  Note that we always have a min step below our trump suit, so we can limit our hands with that response and then make a game try with shapely hands.  This now is true for both majors, not just spades.
	
	General steps are Min, Max, \"Uber (zoom). R+1 is relay for 5-4, 5-5, 6-4, 6+ short, 5 bal, 6 bal (zoom to Q).
	
	\section{+4: Other Major}
	
	\he1--\di2 and \sp1--\he2 both show 5+ cards in the other Major and limit the hand to less than GF values.  Most commonly this is Mixed strength and often will not reach game, but Opener is allowed to come to life with a fit.  Light Invitational hands are possible, but note that \sp1--\he2 is natural, so a very sound invite and only 5 \hhh{s} might choose to start with 1NT instead of \he2.

	We still have \he1--\sp2 and \sp1--\he3 as natural and invitational with 6+, so this is less of a concern with real length in the OM.  
	
	No special bidding here, 2 of either major suggests a contract.  \rem{Perhaps 2NT here should sort of Lebensohl, allowing both a natural 3m suggestion for a contract and a 3m game try in OM.}	
	
	\section{Other}
	
	Inv JS are still on for all suits.  2NT is now the Mixed Raise, with 3M being weak.  This also allows game tries over 2NT if wanted.
	
	Direct double jumps are void splinters, as they were before in TaJ.  
	
	3NT is not currently defined.  It was 17--18 bal before, but that really makes no sense. 1M-TaJ-2x-3NT is choice, so that isn't needed either.  Good 1/4 doesn't much make sense in a limited opener.  Open to suggestions.
	
	\section{Passed Hand}
	
	\rem{This section is more discussion than exposition.}
	
	Everything here is off by a PH.  
	
	Moss/Grue play 4 card majors in 3rd seat with Transfer responses.  It is unclear from the system card whether they start with 1NT always or if \he1--\sp1 is a transfer.  I like the idea and propose we consider doing this; I would have \he1--\sp1 be natural and xfers starting with 1NT, just like we do over a double. R-1 becomes the best raise, so we don't need to strictly have 1M-2M be constructive as we have in the past.
	
	Assuming we do \textit{not} adopt the transfer idea, then everything is natural and NF. 1M-2M would be constructive, pass 1M with a weaker raise.
	  
\end{document}


